\subsection{The arithmetic degree at the selected length}
\label{subsec:deep-arithmetic-degree}

We now estimate the determinant of the component lattice at the selected length.

For the Hilbert--Samuel comparison keep the moving system
of \S\ref{subsec:component-rank-comparison}. Use the two
generating cusp frames at \(u_1,u_2\) and the canonical
section maps at \(v(j)\). The two frames are integral and
their section ratios generate the function field of the flag
component. The \(u_1\) frame gives the hyperplane sections \(x_i\);
the \(u_2\) frame supplies the auxiliary line for the residue and
vanishing tests. At \(j=0,d,T\) the short integral sections span after extension of scalars.
In particular the \(T\)-lattice is full in the determinant
space.

Let \(V\) be the component cusp system at \(u_1\), with integral
lattice \(L\) and Petersson norm at each embedding \(\sigma:F\hookrightarrow\mathbb C\).
A hyperplane block is a vector in \(V\). Choose \(\beta\ge1\) so that
\begin{equation}\label{eq:hs-unit-sphere-beta}
 \sup_\sigma\ \sup_{\substack{X\in V_\sigma\\\|X\|_{\mathrm P,\sigma}=1}}
                      \|X\|_{\sup,\sigma}^2\le\beta .
\end{equation}
For a hyperplane \(X\) and a section \(G\) in a compatible weight,
\[
 \|XG\|_{\mathrm P,\sigma}^2
 \le\|X\|_{\sup,\sigma}^2\|G\|_{\mathrm P,\sigma}^2
 \le\beta\,\|X\|_{\mathrm P,\sigma}^2\|G\|_{\mathrm P,\sigma}^2.
\]
Thus a unit hyperplane costs at most \(\sqrt\beta\) in ordinary length.
An integral generator \(x_i\) of length at most \(\Lambda\) instead
satisfies \(\|x_i\|_{\sup,\sigma}^2\le\beta\Lambda^2\).

Let \(H_Z\) be the normalized Chow height of the moving image in
\(\mathbf P(V^\vee)\), with the lattice \(L\) at finite places and
the Fubini--Study metric induced by the Petersson norm at infinity.
Its Chow form takes \(d+1\) hyperplane blocks, of degree \(D_Z\) each.
A generic line frame \(e\) gives values \(x_i/e\).
Replacing \(e\) by \(a e\) with \(a\ne0\) gives
\[
 [x_0/e:\cdots:x_m/e]=
          [a^{-1}x_0/e:\cdots:a^{-1}x_m/e] .
\]
Thus this common division leaves the projective image and
its normalized Chow height unchanged.

\paragraph{Short generators in the Petersson metric.}
Put \(Z_0=fr\ell\), where \(f\) counts the weight labels,
\(r\) is the rank and \(\ell\) is the logarithmic scale of
the selected model. For a dominant weight \(v\) let
\(S(v)=\sum_{\tau,k}|v_{\tau,k}|\).
The wall gap needed for the short family is
\(8fr+10\), which is available at \(u_1\) and every \(v(j)\)
used here. With coordinates numbered from zero use
\[
         m_0=fr+1,\qquad
         n_{0,\tau}=1+4m_0-v_{\tau,p-1}.
\]
Start at even theta degree \(8m_0\). Adding \(n_{0,\tau}\) makes the last coordinate of the positive block
equal to \(1+4m_0\) before restoring the weight. The wall gap
leaves the required \(4m_0\) interior margin on each
half of the theta exponents.

If \(U\) is the key unit and \(F_{-n_0}\) is the
parallel flat section with shifts \(-n_{0,\tau}\), the
map from the theta family to weight \(v\) is
\[
                  g\longmapsto g\,U^{m_0}F_{-n_0}.
\]
The key unit corrects the central degree. The last
factor subtracts \(n_{0,\tau}\) at each label \(\tau\). Both sections
are nonzero on the component and have nonnegative
orders in its integral rule model. Their flat
identifications over \(F\) carry the full theta span
to the component cusp space at \(v\). The even-degree
theta spanning and norm estimates are applied at the
deep level with the unramified marking used by the
cusp system.

\begin{established}[Short harmonic polynomials at the selected deep level]\label{hyp:deep-short-span}
Take the deep construction with \(N_*=g_0g_1 15^e\) and
\(\ell=L\). The nonterminal lift has positive signature blocks, \([E:\mathbb Q]/2>2\),
good integral models at every prime, and auxiliary level one.
The bounds on the level and the two fields are
\[
 \log N_*\le30\ell,\quad s_{\rm fld}(E)\le3\ell,\quad
 \log(fr)\le\ell,\quad s_{\rm fld}(F)\le391 Z_0^2 .
\]
For the family and the weights used in the short rules, the deep
weight bounds and \(Z_0\le e^{2\ell},\ \kappa\le e^{2\ell}\) give
\begin{equation}\label{eq:deep-short-weight-sizes}
 \begin{aligned}
 S(u_1)&\le Z_0^2(\kappa/T+1)\le e^{7\ell},\\
 S(v(j))&\le Z_0^2\kappa\le e^{6\ell}.
 \end{aligned}
\end{equation}
Both log sizes therefore satisfy the \(10\ell\) bound required by
the short-family estimate.
At \(u_2\) existence
of the auxiliary integral frame uses the ordinary full
lattice result: adding the remainder to \(u_1\) retains
positive consecutive gaps. The short frame entering the
multiplication estimate is at \(u_1\).

Write \(n=8(fr+1)\), \(d_a\) for the harmonic degree and
\[
 K_\theta=d_a+nfr/2.
\]
The even-column stable theta span
uses level classes generating the marked quotient, self-dual
auxiliary lattices \(\Lambda\) locally standard at the primes
of \(Q_S qd N_*\), and compatible auxiliary levels.
The polynomials are in the integral harmonic lattice of
the specified exponent array. They can be chosen with coefficient norm at most \(\exp(C_A Z_0^2 K_\theta)\). The auxiliary dimension satisfies \(n=8(fr+1)\ge8fr+4\);
the entries of the shifted positive block exceed \(n/2\).
Combining the shifts \(n_{0,\tau}\) in the theta weights with \eqref{eq:deep-short-weight-sizes}
and \(\log(fr)\le\ell\) gives the short-rule exponent bound
\[
 K_\theta=d_a+\tfrac12 nfr,\qquad \log K_\theta\le20\ell .
\]
\label{cor:deep-short-readings-1}
\end{established}
\begin{proof}
The level and coefficient-field estimates are those of Lemma~\ref{lem:deep-cover-field-size}; Proposition~\ref{prop:deep-cover-full-lattice} supplies the integral frame at \(u_2\). The powers of \(Z_0\) in \eqref{eq:deep-short-weight-sizes} allow both lists of weights, uniformly in the choice of \(j\). The stable span with a prescribed finite set of standard local lattices is constructed at \ref{cor:deep-cover-denominators-9}. We explain why its harmonic polynomials may be chosen short. Write \(D=d_a\) and \(r=p+q\). The constant polynomial suffices if \(D=0\).

Equip \(\Lambda\) with the real trace norm of its positive Hermitian form. Its trace-dual lattice is \(\mathfrak d_E^{-1}\Lambda\) by Hermitian self-duality. For \(d_E=|\operatorname{disc} E|\), multiplication by \(d_E\) sends this dual into \(\Lambda\), and every nonzero trace norm squared in \(\Lambda\) is a positive integer. Hence Euclidean transference in dimension \(2fn\) gives independent \(v_i\in\Lambda\) spanning over \(\mathbb R\) with
\[
        \lambda_1(\Lambda^*)\ge d_E^{-1},\qquad
        \max_i\|v_i\|\ll fn d_E.
\]
Here either convention for the complex trace norm changes only the absolute constant. \cite{banaszczyk}

For two block partitions \(a^+,a^-\), with trailing entries extended by zero, write the reference joint highest vector as
\[
  P_{a^+,a^-}=\prod_{s=1}^{p}\det(X_{[s],[s]})^{a_s^+-a_{s+1}^+}
          \prod_{t=1}^{q}\det(Y_{[t],[t]})^{a_t^--a_{t+1}^-},
\]
where the reference \(X,Y\) use orthogonal column frames. At each label the row-highest harmonics are generated by block products of leading minors evaluated on pairs of row frames \((A,C)\) with \(h(A,C)=0\). Indeed the mixed contractions of each such product vanish by orthogonality; at the coordinate frame the product is the highest vector of the required joint harmonic representation. Its column-group orbit spans that representation by highest-weight theory, with the determinant exponents equal to the successive differences of the prescribed rows, as in Section~\ref{subsec:theta-row-rank}. In particular, if \(V_a\) is the row-highest harmonic space, the column action satisfies
\[
      V_a=\operatorname{span}\{gP_{a^+,a^-}:g\in\mathrm{GL}_n\}
\]
at each complex label. The allowed coefficient embeddings occur in distinct conjugate pairs; after restriction of scalars their complexifications are independent factors.

Here is an integral parametrization of a dense set of these pairs. Choose row lists \(A,A'\) of length \(p\) and rows \(B_j\) of length \(q\), each a linear combination of the \(v_i\). Set \(H_{ik}=h(A'_k,A_i)\), \(b_{j,i}=h(B_j,A_i)\), taking the form linear in its first entry, and put
\[
 C_j=\det(H)B_j-\sum_k\bigl(\operatorname{adj}(H)b_j\bigr)_k A'_k .
\]
Hermitian integrality puts all these rows in \(\Lambda\) for integral coefficients. The vector of contractions is
\[
   \big(h(C_j,A_i)\big)_i=\det(H)b_j-H\operatorname{adj}(H)b_j=0 .
\] On the open set \(\det H\ne0\) this map onto the orthogonal complement is surjective in the variables \(B_j\), since its restriction there is scalar multiplication. The integral vectors \(v_i\) span the restriction of scalars over \(\mathbb Q\). After base change the parameters can therefore vary independently in each embedding, giving the dense set just used.

The entries of the resulting degree-\(D\) block product are polynomials of degree at most \(d=(2r+1)D\) in the integer parameters. Interpolate in one variable at a time on \(\{0,\ldots,d\}\). Indeed for every linear functional \(\ell_0\) on the coefficient space, successive one-dimensional interpolation gives
\[
   (\ell_0 P)|_{\{0,\ldots,d\}^{\,2fn(2p+q)}}=0
                  \quad\Longrightarrow\quad \ell_0P\equiv0 .
\]
Thus the grid spans.

With \(R=\max(1,2fn d\max_i\|v_i\|)\), the entries of a row in \(A,A',B\) have norm at most \(R\) at every complex embedding. Determinant expansion in the formula for \(C_j\) bounds each entry by \((p+1)!(nR^2)^p R\). Thus the coefficient \(\ell^1\) norm of each block product satisfies
\[
 \log\|P\|_1
 \ll D r\left(\log(2r n)+\log R\right)
 \ll D r\bigl(f\ell+\log(1+D)\bigr)
 \ll Z_0^2 K_\theta.
\]
The second estimate uses \(\log d_E\le 7f\ell\), and the last uses \(Z_0=fr\ell\) and the weight bounds above. The same calculation applies at every embedding of the coefficient field.
\end{proof}

\begin{proposition}[An integral short span on the deep cover]
\label{prop:deep-short-integral-span}
For the deep cover with the data of Lemma~\ref{hyp:deep-short-span},
each even-weight spanning rule \(y\) has an integral
multiple \(sZ_*^2y\), with
\[
     \log\|s Z_*^2 y\|_{\mathrm P}
       \le s_{\rm fld}(F)+O(Z_0^2K_\theta).
\] The determinant-ideal scalar may depend
on the auxiliary lattice. The second scalar depends only on the
level and harmonic degree. The flat-line shift preserves the
rational span with nonnegative order at every integral frame and the
additional archimedean cost in \eqref{eq:parallel-return-norm}.
\end{proposition}

\begin{proof}
\proofstep{Descent of the even-weight span.}
Use the construction at \ref{cor:deep-cover-denominators-9} with
the auxiliary integer divisible by \(Q_S qd\). It realizes the
short span on the identity cover with the standard auxiliary
lattice at every prime dividing \(Q_S qd N_*\). Expand the varying polynomial in short integral harmonic polynomials for that lattice. The resulting theta section depends linearly on the polynomial. The finite class indicator gives the
marked Fourier coefficients, and the orbit of the component
under coefficient embeddings has the prescribed uniformized
Hermitian space. These conditions allow the marked sections and the span of their short theta rules to be compared in the same coefficient space over \(F\).
To state the descent explicitly, let \(V\) be that space and
\(W\subset V\) the span of the specified short rules. If a section \(y\in V\) has complexification in \(W\otimes_F\mathbb C\),
then
\[
 (y+W)\otimes1=0\ \text{in }(V/W)\otimes_F\mathbb C
                 \quad\Longrightarrow\quad y\in W .
\]

\proofstep{The determinant-ideal scalar.}
Two different scalars enter. One depends on \(\Lambda\).
Let \(I\) be the index of its intersection with the
standard lattice in \(\Lambda\), and \(\mathfrak a\) the
index ideal, of norm \(I\). In each embedded determinant
line the intersection contributes \(\tau\mathfrak a\).
Take \(B=\prod_\tau(\tau\mathfrak a\,\mathcal O_F)^{m_\tau}\)
with the signature multiplicities. Its norm is
\(I^{r[F:\mathbb Q]/2}\) since each conjugate pair
has total multiplicity \(r\).

In the totally complex
coefficient field, apply the convex body theorem to
\(B^{-1}\) with disk radii
\(\exp(s_{\rm fld}(F)) I^{-r/2}\). The lattice
covolume is \(2^{-[F:\mathbb Q]/2}
\sqrt{|\operatorname{disc} F|}/\mathrm N B\).
The disk product exceeds \(2^{[F:\mathbb Q]}\) times this
covolume, by the field-size bound and \(\pi>2\).
There results \(s\ne0\) in \(B^{-1}\) with
\[
   |\sigma s|^2 I^r\le \exp(2s_{\rm fld}(F))
           \qquad\text{for every }\sigma.
\]
Membership is precisely the determinant-ideal inclusion
used for the integral short rule. The same \(s\) pays
for the factor \(I^{r/2}\) in its analytic bound.

\proofstep{A second scalar for finite prime and level denominators.}
The other scalar can be fixed before the auxiliary lattice
and the class tuple are chosen. Set
\[
    Z_* = N_*^{frn+d_a}\,|\operatorname{disc}E|^{d_a}.
\]
The level class average is cleared by its cardinality
\(N_*^{frn}\), and the polynomial and eigenframe factors
by
\[
 N_*^{d_a}|\operatorname{disc}E|^{d_a}.
\]

In the finite order test write the raw rule \(y\) as
\(\beta_\Lambda\) times the short rule in determinant
coordinates. The class indicator at each label has integral
Fourier coefficient (a root of unity). With \(s\) in the inverse index ideal, the coprime torsion comparison gives 
\[
 v(sZ_*y)\ge0\qquad(v\nmid N_*).
\]
 At a prime dividing \(N_*\) the auxiliary
lattice is locally standard, so \(s\) is integral there.
Use the chart comparison at that prime for the short
rule in determinant coordinates, and multiply back by
\(\beta_\Lambda\).

For a prime dividing the level once
the class-sum bound clears this with \(Z_*\); for one
whose square divides the level the trivial-lattice
comparison clears it with \(Z_*^2\). In the latter comparison the class bound is applied after a finite extension. For an element \(a\) of the base and integer-normalized orders upstairs and downstairs, the normalization is
\[
 v_{\rm base}(a)=e^{-1}v_{\rm ext}(a),
\]
where \(e\) is the ramification index. Summing integral Fourier coefficients
preserves the lower order and descends it to the original
test. Combining the local tests gives
\[
 v(sZ_*^2y)\ge0\qquad(\text{every finite test place}),
\]
including the two primes of the growing principal depth. Moreover
\[
  \log Z_*^2
    =2(frn+d_a)\log N_*+2d_a\log|\operatorname{disc}E|
    \ll Z_0^2 K_\theta
\]
since \(frn\le2K_\theta\) and
\(\log|\operatorname{disc}E|\le7f\ell\).
\label{cor:deep-short-readings-2}

\proofstep{The common complex norm bound.}
For the complex comparison choose, at an embedding with
positive supremum, an orthonormal flag frame reaching
at least half the supremum and with squared first minimum
at least \(e^{-80\ell}\). This is the level point
estimate for spread at most \(3K_\theta\): its level
cutoff is at least \(e^{-57\ell}\) and the loss in
squared first minimum is at most \(\operatorname{rd}(E)^2\).
The Fourier transform of the point indicator has absolute
value one at each label.

On that frame the classical
self-dual theta norm estimate with this first minimum,
the polynomial coefficient bound and its determinant factor
gives
\[
   \|sy\|_{\sup}\le
   \exp(s_{\rm fld}(F))\,
   \exp\big((C_A+C_B)Z_0^2 K_\theta\big).
\]
The common Lie factor in this estimate is
\[
 (2\pi\xi_R)^{-d_a/2};
\]
the index factor is paid
by the preceding inequality for \(s\). The uniformization
uses the probability mass, so the same bound holds for
the Petersson norm on each component. Multiplication by
\(Z_*^2\) and \(K_\theta\ge1\) absorb the remaining
field and integer costs. All these scalar multiples are
nonzero multiples of the respective spanning rules.
Restriction to \(Z\) therefore gives an integral
Petersson-bounded span at the even theta weight.
\label{cor:deep-short-readings-3}

\paragraph{Making the flat weight return integral.}
The key unit and the parallel return are used with
different estimates. For the parallel return let
\(r_\sigma>0\) be its constant complex norms and
\(\rho_w\) its finite orders before scalar adjustment.
They have finite order support, conjugate embeddings
have the same norm, and outside the bad rational primes
the orders are integers by the unramified frame test.
In \(F\) form the fractional ideal
\[
      J=(b_{\rm bad})\prod_w\mathfrak p_w^{\lfloor-\rho_w\rfloor},
             \qquad b_{\rm bad}=\prod_{l\,\mathrm{bad}} l .
\]
Every nonzero \(a\in J\) satisfies
\(v_w(a)+\rho_w\ge0\). Off the bad primes this follows
by integrality; at a bad prime \(v_w(b_{\rm bad})\ge1\)
covers the rounding.

If \(h=[F:\mathbb Q]\), then
\[
 \tfrac1h\log\mathrm N J\le\log b_{\rm bad}
            -\tfrac1h\sum_w\rho_w\log\mathrm N w.
\]
Use a symmetric Minkowski body with radii
\(A/r_\sigma\), where
\(A=(5/2)(\mathrm N J\sqrt{|\operatorname{disc}F|}
\prod_\sigma r_\sigma)^{1/h}\). Its volume exceeds
the Minkowski threshold for \(J\). The resulting nonzero
element gives a single scalar at all embeddings. Taking
logs yields
\begin{equation}\label{eq:parallel-return-norm}
  \log(|\sigma a|r_\sigma)\le
  -\deg_{\rm flat}+\tfrac12\log\operatorname{rd}(F)
             +\log b_{\rm bad}+1 .
\end{equation}
Here the degree uses the archimedean negative log norms
and the finite additive orders. The preceding flat
calculation bounds the negative degree by the expression
linear in the absolute exponents below. Thus the
parallel twist has nonnegative order at every integral
frame and a simultaneous upper norm with this loss. It
is invertible on the generic component, so multiplication
preserves the span. \label{cor:deep-short-readings-4}
\end{proof}


For each member the stable theta norm costs at most
\(C_g Z_0^2 K_{\theta}\) in logarithmic length, where
\(K_{\theta}\) is the absolute sum of the theta weight
coordinates. The shift of each coordinate by \(n_{0,\tau}\)
and the theta degree gives
\(K_{\theta}=O(rS(v)+fr(fr+1))\). The sum of the absolute
flat exponents is bounded by \(S(v)+f(1+4m_0)\).
The key unit costs \(2fr m_0\log(2\pi)\).

Write \(\operatorname{rd}\) for the root discriminant, and
let \(c_1\) be the logarithmic flat period.
The flat
section costs at most
\[
 (|c_1|+q\,\log\operatorname{rd}(E))
     \sum_\tau |n_{0,\tau}|
 +\sum_{l\ {\rm bad}}\log l
 +\tfrac12\log\operatorname{rd}(F)+1 .
\]
The selected
bounds are \(|c_1|\le9r\ell\),
\(\log\operatorname{rd}(E)\le3\ell\) and
\(\log[F:\Q],\log\operatorname{rd}(F)\le391Z_0^2\).
The bad-prime sum on this level is \(O(Z_0^2)\).
Together these give a spanning family of integral
vectors of log length at most
\(C(Z_0^3S(v)+Z_0^5)\) with an absolute constant \(C\).
\label{cor:covolume-hilbert-samuel-1}

\begin{lemma}[Short integral generators at the selected weights]
\label{lem:deep-integral-generators}
On the component of Lemma~\ref{hyp:deep-short-span}, use
the integral span of Proposition~\ref{prop:deep-short-integral-span}
over the same coefficient field as the determinant. For
\(\kappa\ge1\), \(P\ge Z_0^{12}\) and sufficiently large \(Z_0\),
let \(h(\lambda)\) be the largest log Petersson norm among the chosen generators of the full integral lattice at \(\lambda\). One can choose them so that
\[
 \begin{aligned}h(u_1)&\le P(\kappa/T+1),& h(v(j))&\le P\kappa\quad(0\le j\le T),\\h(v(0))&\le dP(\kappa/T+1).\end{aligned}
\]

\end{lemma}

\begin{proof}
\proofstep{Passing from a short span to integral generators.}
To obtain generators of the integral lattice, choose
an \(F\)-basis from this short span. If \(h=[F:\Q]\)
and its size is \(a_0\), tensor it with \(h\) independent
integers of \(F\) of conjugate absolute values at most
\(h^3\operatorname{rd}(F)\). Such integers follow from
the successive-minima bound for a number field.
This is a \(\Q\)-basis contained in the integral lattice.
Call the \(ha_0\) vectors of the resulting basis \(v_i\).

For each vector \(z=\sum_i c_iv_i\) in an integral basis of the lattice,
retain the remainder
\[
 z'=z-\sum_i\lfloor c_i\rfloor v_i
   =\sum_i(c_i-\lfloor c_i\rfloor)v_i,\qquad
 \|z'\|\le ha_0\max_i\|v_i\|.
\]
The \(v_i\) and the remainders together recover every original \(z\). The new list of at most \(2ha_0\)
vectors generates the lattice and increases the length
bound by at most \(\max(1,h^4a_0\operatorname{rd}(F))\).
This also proves the claim when the chosen \(F\)-basis
generates a sublattice of large index.

The selected size hypotheses also give
\[
 \log\max(1,S(v))\le20\ell,
\]
the range of the uniform
point estimate. The point evaluation bound in the normalized Petersson
model with summed spread at most
\(1+2[E:\Q]S(v)\) gives a Bergman bound with logarithm
at most \(4000Z_0^3\). Its integral over a component,
with its mass bound, controls the log rank by
\(O(Z_0^3)\). The resulting generator count has log
bounded by \(O(Z_0^3)\), and the cost of passing from
spanning vectors to integral generators is absorbed in
\(C(Z_0^3S(v)+Z_0^5)\). At \(u_1\) this is precisely the unit-ball estimate \eqref{eq:hs-unit-sphere-beta}. Thus
one may take \(\log\beta\le4000Z_0^3\le P\) after increasing the fixed
threshold on \(Z_0\). The lengths of the integral generators enter separately.

\proofstep{The size and norm at the actual weights.}
The weight estimates are
\[
 S(u_1)\le Z_0^2(\kappa/T+1),
\]

\(S(v(j))\le Z_0^2\kappa\) for \(0\le j\le T\), and
\(S(v(0))\le 3d Z_0^2(\kappa/T+1)\).
The last bound follows by summing the three terms
\(\kappa_{\rm can},u_2,d u_1\) at \(j=0\).
With \(\kappa\ge1\), \(P\ge Z_0^{12}\), and \(Z_0\)
eventually larger than the fixed constants,
the log generator bounds at \(u_1\) and \(v(j)\)
are respectively \(P(\kappa/T+1)\) and \(P\kappa\).
At \(v(0)\) the separate bound is \(dP(\kappa/T+1)\).
For example the latter cost is at most
\[
 4Cd Z_0^5(\kappa/T+1)\le dP(\kappa/T+1).
\]

The log rank bound at \(v(T)\), combined with
\(D_Z\le J_Z\le r_Z\), gives \(\log D_Z\le P\).
All these are estimates on the same component with its
prescribed determinant frame. \label{cor:covolume-hilbert-samuel-2}

\end{proof}

\paragraph{Chow form and determinant.}

The short sections determine a projective frame. The height term in the arithmetic degree is obtained by comparing its Chow form with a determinant on the same component.

The determinant estimate uses a short generating family
of the \(u_1\) frame, with number \(n_1\) and length bound
\(\Lambda\). It uses ordinary length bounds
\(\Lambda_0,\Lambda_d,\Lambda_T\) at \(0,d,T\). Put
\(I_0=\log\Lambda_0\le dP(\kappa/T+1)\).
The length and squared norm bounds are at least one.
Lemma~\ref{lem:deep-integral-generators} supplies the following
bounds on the selected component
\begin{align*}
 &\log n_1,\ \,\log\beta,\ \,\log D_Z\le P,
 \qquad \log\Lambda\le P(\kappa/T+1),\\
 &\log\Lambda_d,\ \,\log\Lambda_T\le P\kappa,\\
 &I_0\le4\{dP(\kappa/T+1)+TP\}+2dP\kappa/T.
\end{align*}
The radius and the volume correction use
\[
 \log R=\log\Lambda+26\log\big(n_1 d D_Z
                 (2+|H_Z|+\log\Lambda)\big),
 \qquad
 g=\frac{(d+1)(\log n_1+1)}{2}+1 .
\]

\begin{proposition}[The selected-length Hilbert--Samuel lower bound]
\label{prop:component-arithmetic-degree}
For the moving system at the selected length \(T\), the normalized
ordinary degree of the component lattice satisfies
\eqref{eq:deep-hs-raw}; in particular, writing \(d_{\rm cusp}\) for this normalized degree, \eqref{eq:deep-hs-budget} gives
\[
 d_{\rm cusp}\ge J_Z\{T H_Z-3(d+1)|H_Z|-A\log(2+|H_Z|)-R_0\},\qquad R_0=4000b+2c_{\rm flat}.
\]

Both estimates use the same \(D_Z\) and \(J_Z\) as the
component rank calculation.
\end{proposition}

\begin{proof}
\proofstep{The determinant of the integral moving system.}
Here is the determinant estimate on this integral
system. Use generic \(d\) hyperplane cuts of the
moving system. On the open set where the cut is reduced
and meets the flag locus, its residue algebra has rank
\(D_Z\). After a splitting extension the specialized
Chow form in the last hyperplane block is, up to a
nonzero scalar, the product of the \(D_Z\) linear
evaluation forms. The filtration by powers of the cut
ideal has monomial symbols in its \(d\) generators.
The residue-surjectivity test uses the second frame
and the nef and big adjoint lines. The cohomology
vanishing test identifies the last filtration piece
with the products from degree \(d\).

In particular
the lower symbol columns have \(D_Z\) residue columns
for each multi-index of total degree less than \(T-d\).
Counting weak compositions in \(d\) variables gives the number of lower columns,
\[
 J_Z=D_Z\sum_{j=0}^{T-d-1}\binom{j+d-1}{d-1}=D_Z\binom{T-1}{d}.
\]
Take a cut maximizing
the deepest filtration rank \(k\) on this open set. Complete
the symbol columns first from the degree-\(d\) family
and then from the degree-\(T\) family, using the full
spans at these degrees. The result is a nonzero
polynomial determinant with columns integral over the
coefficient lattice. Over every good specialization
the symbol of order \(|\alpha|\) has \(T-|\alpha|\)
factors of the last hyperplane. Its determinant over the
residue algebra gives the specialized Chow form to that
power.

Summing over \(|\alpha|<T-d\) gives the exponent
\[
 T\binom{T-1}{d}-d\binom{T-1}{d+1}
   =\binom{T}{d+1}+d\binom{T-1}{d}.
\]
The generic rank choice bounds every specialized last
piece by the number of chosen degree-\(d\) columns.
The Chow polynomial is irreducible and has positive degree
in the last block, so it is primitive over the polynomial
ring in the first \(d\) blocks. Gauss's lemma descends the
divisibility over the fraction field of that ring to the
original determinant. This uses the dimension of the original
residue algebra, which is \(D_Z\) before as well as
after splitting.
\label{cor:covolume-hilbert-samuel-3}

\proofstep{The Chow form on a component.}
Write \(f_Z\) for the Chow form and factor out
\(f_Z^{\binom{T}{d+1}}\), which suffices for the degree bound.
Write \(Q\) for the resulting cofactor.
The three column counts are
\[
 J_Z,k,r_Z-J_Z-k.
\]

Their degrees in the hyperplane variables, summed over
blocks, are \(T,T-d,0\), respectively. The extracted
power \(f_Z^{\binom{T}{d+1}}\) has total degree
\[
 (d+1)D_Z\binom{T}{d+1}=TJ_Z.
\]
Therefore the total
degree of \(Q\) is \(k(T-d)\).

For a scalar multiform \(B\) on \(V^{d+1}\), use the mean of
\(\log|B|\) over the product of Petersson unit spheres at infinity.
At a finite prime \(\mathfrak p\), write
\[
 q_{\mathfrak p}:\mathcal O_{F,\mathfrak p}^{\,n_1}\twoheadrightarrow L_{\mathfrak p},
 \qquad (z_i)_i\longmapsto\sum_i z_i x_i .
\]
Use the maximum absolute coefficient of \(B\circ(q_{\mathfrak p},\ldots,q_{\mathfrak p})\).
This agrees with its Gauss norm in a local lattice basis. Each set of generators
is an integral linear combination of the other, and neither substitution
increases the maximum. The normalized sum of the finite logarithms
and the infinite means is \(h_{\rm Mahler}(B)\).

Let \(n_{\rm im}\le n_1\) be the dimension of the
hyperplane coefficient space before this substitution.
The geometric Chow height uses the same finite Gauss
norms and adds the harmonic correction per homogeneous
degree to the sphere mean. Dividing by \((d+1)D_Z\) gives
\[
 h_*:=\frac{h_{\rm Mahler}(f_Z)}{(d+1)D_Z}
 = H_Z-\frac12\sum_{a=1}^{n_{\rm im}-1}a^{-1}
 \ge H_Z-\frac{\log n_1+1}{2}.
\]
The lower bound for a nonzero multiform of total degree \(a\)
in this frame is \(-c_Ha\), where
\[
 c_H=\log n_1+\log\Lambda+(\log n_1+1)/2.
\]
Here the two norm comparisons use the full integral frame,
its spanning evaluations and the bound \(\Lambda\).
The unit-sphere log integral has a loss at most
\((\log n_1+1)/2\) per homogeneous degree.

To see the
sphere bound, rotate a point of maximum modulus to the
first coordinate axis and average in the phase of that
coordinate, using Jensen's inequality on the polynomial
in that phase. The leading coefficient is the value at
the first axis multiplied by the corresponding power of
the coordinate modulus. For a uniform point \(z\) on the unit
sphere in \(\mathbb C^{n_{\rm im}}\) the mean of \(\log|z_1|\) is
\[
 -\frac12\sum_{a=1}^{n_{\rm im}-1}a^{-1},
\]
at least
\(-(\log n_1+1)/2\). Apply the bound in each block.

For the cofactor lower bound choose any nonzero coefficient
after the spanning-family substitution. Such a coefficient
exists because the family spans over the field. The finite
Gauss norms bound the coefficient locally. At infinity
Cauchy's formula on the coordinate unit torus bounds it by
the supremum over linear combinations of the generators
with coefficients of modulus one. Each combination has
frame norm at most \(n_1\Lambda\). Homogeneity and the
sphere estimate bound that supremum by the exponential
of the sphere mean plus \(c_H a\). The product formula
for the coefficient gives
\[
 h_{\rm Mahler}(Q)\ge-c_Ha.
\]


\proofstep{Multiplicativity of the height.}
Mahler height is additive on products. Thus the determinant
has height at least \(TJ_Zh_*-c_H k(T-d)\).
Let \(L_T\) be the full lattice at \(v(T)\) and choose an \(F\)-basis
\(\xi\) of its determinant line. Write
\[
  \det L_T=\mathfrak a\,\xi,\qquad
  \bigwedge_{i=1}^{r_Z} C_i(X)=\delta(X)\,\xi ,
\]
where \(C_i(X)\) are the three families of polynomial columns just constructed.
At every finite prime each column coefficient after the substitution
\(q_{\mathfrak p}\) is a sum of integral products of sections and lies in
\(L_{T,\mathfrak p}\). Multilinearity puts each coefficient of the wedge
in \(\mathfrak a_{\mathfrak p}\xi\), so the finite Gauss norm of \(\delta\)
is at most the absolute value of a local generator of \(\mathfrak a\).

At an infinite place evaluate this same wedge on
\(\|X_0\|_{\mathrm P}=\cdots=\|X_d\|_{\mathrm P}=1\).
A column starting at degree \(j=0,d,T\) is a product of \(T-j\)
hyperplanes and a section of length at most \(\Lambda_j\).
Repeated multiplication gives
\[
 \|C_i(X)\|_{\mathrm P}\le\beta^{(T-j_i)/2}\Lambda_{j_i},
 \qquad j_i\in\{0,d,T\}.
\]
Here \(j_i\) is the starting degree of column \(i\). Hadamard bounds
\(\log|\delta(X)|+\log\|\xi\|_\sigma\) by the sum of the log column
bounds, with multiplicities \(J_Z,k,r_Z-J_Z-k\). Average on the spheres
and add the finite bounds. The terms containing the fractional ideal
and the basis vector combine to give
\[
 \widehat{\deg}_Z(v(T))=
 \frac{-\log N_{F/\mathbb Q}\mathfrak a-\sum_{\sigma:F\hookrightarrow\mathbb C}
                \log\|\xi\|_\sigma}{[F:\mathbb Q]} .
\]
Combining this upper estimate for \(h_{\rm Mahler}(\delta)\) with its
Chow-factor lower estimate gives
\begin{align*}
\widehat{\deg}_Z(v(T))\ \ge\
 &J_Z\bigl(T h_*-I_0-\tfrac T2\log\beta\bigr)\\
 &-k\bigl((T-d)c_H+\log\Lambda_d
                        +\tfrac{T-d}{2}\log\beta\bigr)\\
 &-(r_Z-J_Z-k)\log\Lambda_T .
\end{align*}

\proofstep{The contribution per lower column.}
For each of the \(J_Z\) lower columns the ordinary length loss from
the hyperplanes is at most
\[
             \tfrac T2\log\beta\le TP/2\le\sqrt b/2.
\]
The factor \(\log\Lambda\) from substituting the integral generators
in the cofactor bound has coefficient \(k(T-d)\). Using the rank
share \(k\le r_Z-J_Z\le2zJ_Z\), we get
\[
 \frac{k(T-d)}{J_Z}\log\Lambda
 \le 2zTP(\kappa/T+1)\le 3b
\]
by \(zP\kappa\le b,\ z\le1/2,\ TP\le\sqrt b\le b\).
Since \(c_H\le\log R\), the complementary bracket below
dominates the costs per column in both remaining families.
The loss \(T(\log n_1+1)/2\) on the first line is at most
\((T+d^2)(g+2)\). Adding the other nonnegative error bounds gives
\begin{align}
 \widehat{\deg}_Z(v(T))
 &\ge J_Z\big\{T H_Z-3(d+1)|H_Z|-I_0
       -26d^2(\log R+\log T)\notag\\
 &\hspace{41mm}{}-(T+d+1)\log\beta
       -(T+d^2)(g+2)\big\}\notag\\
 &\quad -(r_Z-J_Z)
       \big\{T(\log 2+\log R+\tfrac12\log\beta)
                        +\log\Lambda_d+\log\Lambda_T\big\}.
 \label{eq:deep-hs-raw}
\end{align}
The last bracket is
nonnegative and controls the entire complement to the
symbol columns. The degree is in the ordinary norm
convention, before passing to the squared-norm slope.
\label{cor:covolume-hilbert-samuel-4}

\paragraph{Bounding the error on one component.}
Write \(X=\log(2+|H_Z|)\) and \(Q=\kappa/T\).
The parameter bounds include \(1\le P\), \(1\le b\),
\(\log b\le P\), \(TP\le\sqrt b\) and
\(zP\kappa\le b\). From the radius formula and
\(\log(1+\log\Lambda)\le\log\Lambda\) we obtain
\[
 0\le\log R\le 26X+27P(Q+1)+65P .
\]
Here \(\log d\le\log T\le P/2\) and
\(2+|H_Z|+\log\Lambda\le(2+|H_Z|)(1+\log\Lambda)\).
We have \(r_Z\le2J_Z\) and
\(r_Z-J_Z\le 2zJ_Z\) from the rank shares.
Thus the coefficient of \(X\), after factoring out
\(J_Z\), is at most
\[
 A=26(26d^2+2zT).
\]

\proofstep{Uniform bounds for the other terms.}
For the other terms use
\(d^2P Q\le b/7\) by \(7d^2/T\le z\),
\(d^2P\le b/14\) by \(z\le1/2\), and
\(TdP\le(TP)^2\le b\).
In particular \(I_0\le9b\) and the radius term
with coefficient \(26d^2\), after removing \(A\)'s
share of \(X\), costs at most \(300b\).
The two main terms with \(\log\beta\) and \(g+2\)
cost at most \(2b\) and \(10b\) respectively.
For the complementary columns use
\[
 zTPQ=zP\kappa \le b,
\]
\(zTP\le b/2\), and \(\log2\le1\le P\).
The sum of their constant terms is at most \(200b\).
These error terms are bounded by \(4000b\).

\proofstep{The degree loss from the flat line.}
Let \(c_{\rm flat}\ge0\) bound the degree lost per unit of rank
when twisting by the flat determinant line to pass from weight \(v(T)\)
to weight \(u_M\) on this component.
The geometric cusp ranks at \(v(T)\) and \(u_M\) agree by
the flat weight comparison.
Twisting back decreases the degree by at most \(r_Z c_{\rm flat}\le
2J_Z c_{\rm flat}\). Thus the lower ordinary degree on the
deep geometric cusp lattice at \(u_M=u_W-c_0\omega_{\rm H}\) is
\begin{equation}
  J_Z\big\{T H_Z-3(d+1)|H_Z|
       -A\log(2+|H_Z|)-R_0\big\},
  \qquad R_0=4000b+2c_{\rm flat}.
  \label{eq:deep-hs-budget}
\end{equation}
In this calculation \(A\le702T^2\),
\(\log\max(1,A)\le8P\) and \(T^2P\le b\),
so
\[
 A\log\max(1,A)\le6000b.
\]

In this twist from \(v(T)\) to \(u_M\), the sum of the absolute
values of the exponents is at most \(6fr\).
In each component the common flat period and the
relative different cost at most
\((9r\ell+3q\ell)6fr\). The bad-prime and discriminant
allowance for this shift is at most \(1+6\ell+6f\ell+30\ell\).
These are bounded together by \(120 Z_0^2\).
Thus \(c_{\rm flat}\le C_p Z_0^2\) with \(C_p=120\),
and \(Z_0^2\le b\) on the chosen scales.
Consequently the cost \(A\log\max(1,A)+R_0\)
is at most \((10000+2C_p)b\), with the same \(D_Z\)
and \(J_Z\) as in the rank calculation.
\label{cor:covolume-hilbert-samuel-5}
\end{proof}
